\section{General Discussion}
\label{sec:discussion}
\subsection{Practical implications for CPU encryption systems}
The measurements indicate that cryptographic-library selection and operation design should be considered jointly. OpenSSL's substantial transform advantage on Windows becomes smaller when the timer includes the surrounding instrumented operation. The fastest isolated cipher implementation therefore need not deliver a proportionally faster complete request, and a different request design could change the ranking. The observed ranking applies to this benchmark, not a production service.

Selecting an optimized native library matters most when encryption processing occupies a substantial part of the relevant operation. For bulk workloads, the Arm64 contrast makes capability selection an important configuration check. At broader interfaces, allocation, context initialization, conversion and result handling also deserve attention. Engineers should measure their actual interface and payload distribution, then test whether persistent contexts, reusable output buffers or fewer conversions reduce its cost. These remain testable opportunities, not demonstrated improvements.

Short authenticated messages require separate evaluation. The GCM encryption and verified-decryption rates at 1 KiB differ markedly from those at 1 MiB and above, even though cipher-object construction is excluded throughout. Large-buffer throughput consequently cannot characterize complete short-message requests. An application-level evaluation should include the required setup, AAD and tag operations, as well as the surrounding request processing.

Cloud users should distinguish advertised instruction availability from library configuration and actual execution. A practical audit records the loaded cryptographic library, version, provider and capability overrides, then checks capability settings in that same runtime. OpenSSL exposes CPU settings through its diagnostic interface~\cite{opensslinfo}; a separate command-line installation should not be assumed to represent the library loaded by an application. Correctness checks and interface-specific measurements can reveal slow configurations; instruction tracing is needed to establish the executed path.

\subsection{Limitations and scope of generalization}
Different processors, operating systems, libraries and runtimes prevent a controlled x86-versus-Arm comparison. Five sessions per configuration are repeated observations on one system, not five independent systems. Scheduling, frequency and cloud contention were not fully controlled, and prepared-buffer reuse differs from cold-cache or streaming workloads. The Oracle worker campaigns are nonpaired and have different aggregate byte counts.

The timing scopes also constrain interpretation. OpenSSL's native transform and PyCryptodome's Python-level transform include different interface costs. The broader operation includes benchmark diagnostics and bookkeeping, so neither throughput retention nor the change in library ratio measures a production request or identifies an individual overhead component. \texttt{OPENSSL\_armcap=0} changes several capabilities at once, and no instruction-level execution trace was collected. The observed effects of implementation and timing scope are empirical findings; explanations in terms of particular instructions, allocation costs or conversion costs remain hypotheses requiring controlled profiling.

CBC, GCM and CTR originate from separate campaigns with different timing conventions, preventing a matched estimate of authentication cost. Functional validation covers the tested outputs and GCM rejection cases, not side-channel resistance. No complete storage or messaging application, deployed key-management workflow or energy profile was measured.

The strongest generalization concerns experimental design: evaluate the required functionality at the relevant interface and distinguish processing from surrounding-operation costs. Numerical ratios remain specific to the recorded configurations.

\subsection{Applications and analogies in other domains}
Bulk storage and messaging services should test implementation advantages after preparation and result delivery, using their actual payload distributions. The Oracle observation also motivates measuring aggregate completion under concurrency separately from single-request latency.

The same measurement issue arises in compression, numerical computing and machine-learning inference. Buffer preparation can surround a fast compression routine, data conversion can precede an optimized numerical kernel, and preprocessing can accompany inference. These are conceptual analogies: a component-level advantage must be evaluated within the larger operation that uses it. No measurements from those domains are claimed here.

For reproducibility, a report should specify the cipher mode, key size, payload distribution, nonce and AAD policy, loaded library and capability configuration, buffer ownership and exact timer placement. It should distinguish per-worker from aggregate throughput, identify the byte count used at each boundary, and preserve session-level observations rather than only the best rate. These details make it possible to reproduce a contrast and assess its relevance to another workload.

\subsection{Future research}
The first priority is a comparison with matched native entry boundaries for OpenSSL and PyCryptodome, with controlled alternatives for context and output-buffer reuse. This would separate implementation differences from language-interface and preparation costs. Instruction-level profiling on Arm64, combined with selective validation of cryptographic capability settings where supported, would provide stronger evidence about the optimized paths responsible for the observed contrast.

A second priority is equivalent AES-GCM encryption and verified-decryption workloads on Arm64, preserving nonce, AAD and tag policies and measuring setup separately. Complete request-level benchmarks should then use realistic message-size distributions, concurrency and input/output costs to test whether the observed advantages survive within a service. If suitable instrumentation becomes available, energy per processed byte and per completed request would provide an additional criterion alongside throughput and latency. These studies remain future work.
